\section{The odd witnesses also kill every abelian-cover homology action}

Let \(F=\widehat F(x,y)\), and let \(M\) be the multiplier-one
subgroup of full profinite GT.  Use the normalized full adelic
Alexander unit \(B_g(X,Y)\), with its semilinear cocycle law
\[
 B_{gh}(X,Y)=B_g(X,Y)B_h(X^{\lambda_g},Y^{\lambda_g}).
\]
This is the full completed group-ring coefficient, not only its
augmentation expansion at one coefficient prime.  The convention
is \cite[\S2.3, p.~101]{Furusho2003}.  On multiplier one it is
the based coefficient on the Alexander module.

\begin{lemma}\label{lem:full-integral-alexander-unitarity}
For every \(g\in M\),
\begin{equation}\label{eq:full-alexander-unitarity}
                    B_g(X,Y)B_g(X^{-1},Y^{-1})=1
\end{equation}
in the full completed adelic group algebra.  Consequently,
for complex conjugation \(c\) and any \(z\in M\),
\begin{equation}\label{eq:full-alexander-symmetrization}
                           B_{z c z c^{-1}}=1.
\end{equation}
\end{lemma}

\begin{proof}
Fix a positive integer \(d\), a coefficient prime \(r\), and a
character \((a,b)\) of \((\mathbf Z/d)^2\) with values in a
finite extension \(L/\mathbf Q_r\).  The cover corresponding to
\[
                         F_d=\ker(F\to(\mathbf Z/d)^2)
\]
is the punctured Fermat cover.  The automorphism of \(F\)
associated with \(g\) preserves \(F_d\), acts trivially on its
deck group, and preserves every positive peripheral class.
For every nontrivial deck character its eigenspace in
\(H_1(F_d,L)\) has dimension one, and its scalar action is
\(B_g(a,b)\).  This follows either by specializing the
Alexander module or by the cellular complex with differential
\((a-1,b-1):L^2\to L\).

We specify the boundary characters and their scalars.  In
the usual normalization \(x\) is fixed and \(y\) is conjugated
by an element of \(F'\).  Write \(z_0=(xy)^{-1}\) and
\(\phi_g(z_0)=k^{-1}z_0k\).  If \([k]=(u,v)\) in
\(F^{\rm ab}\), then \(u=v\): modulo \(\gamma_3F\),
the automorphism is the identity, so the commutator of
\([k]\) with \((-1,-1)\) is zero.  The degree-two free
commutator coordinate gives \(u-v=0\) in
\(\widehat{\mathbf Z}\).

The boundary subspace is the sum of the permutation modules
induced from the three cyclic peripheral subgroups, with one
trivial relation.  Its nontrivial characters satisfy
\(a=1\), \(b=1\), or \(ab=1\).
The conjugator at \(x\) is trivial, the one at \(y\) has
trivial deck image, and the one at \(z_0\) has diagonal deck
image.  Their scalars on the respective boundary characters
are therefore one; in the last case it is
\(a^u b^v=(ab)^u=1\).
Hence \(B_g(a,b)=1\) at every boundary character.
At \((a,b)=(1,1)\), the scalar is also one because the
class-two commutator is fixed.

The homology modulo boundary classes has the nondegenerate
alternating form of the compactified cover.  The induced
automorphism preserves this form with multiplier one.
Algebraically, in a surface presentation, preservation of
each positive boundary word preserves the degree-two
fundamental relation \(\sum_j[A_j,B_j]\) after quotienting
by the boundary classes.  Conjugating a boundary word adds
only wedges involving its boundary class, which are zero in
that quotient.  This gives preservation of the associated
alternating form without asserting that the profinite
automorphism is realized by a discrete homeomorphism.

The deck action preserves the same form.  Thus the
\((a,b)\) and \((a^{-1},b^{-1})\) eigenspaces pair
nondegenerately whenever \(a,b,ab\ne1\), and
\[
                         B_g(a,b)B_g(a^{-1},b^{-1})=1.
\]
The boundary calculation proves this identity in all
remaining cases as well.  It is valid also when \(r\mid d\),
since the character decomposition is being taken over the
characteristic-zero field \(L\).

For a fixed finite deck group, its characters over a splitting
field separate the characteristic-zero group algebra.
The integral algebra \(\mathbf Z_r[(\mathbf Z/d)^2]\)
embeds in that algebra.  The identity therefore holds in
this integral algebra for every \(r,d\).  Passing to the
inverse limit proves \eqref{eq:full-alexander-unitarity}
in the full adelic completed algebra.

Because this algebra is commutative, the cocycle law gives
\[
 B_{a z a^{-1}}(X,Y)=B_z(X^{\lambda_a},Y^{\lambda_a})
                                            \qquad(z\in M).
\]
Take \(a=c\), whose full multiplier is \(-1\).
Equation \eqref{eq:full-alexander-unitarity}, followed by
the multiplier-one product law, proves
\eqref{eq:full-alexander-symmetrization}.
\end{proof}

\begin{corollary}\label{cor:odd-target-invisible-abelian-covers}
Let \(J\) be the common radical, primewise-word, and signed-CM
kernel from the preceding theorem, and put
\[
                              J_B=J\cap\ker B.
\]
For every odd prime \(p\), restriction
\[
                  Y_p\longrightarrow H^1(J_B\cap H,\mathbf F_p)
\]
is injective.  Every witness \(w_{ij}\) from that theorem
lies in \(J_B\cap H\), with the same diagonal evaluation
\(\kappa_{ab}(w_{ij})=2\delta_{ai}\delta_{bj}\).
Each witness acts trivially on the full adelic first
homology of every finite abelian cover of the punctured line.
In particular, Fermat-cover homology alone cannot provide
a full GT lift of a nonzero odd target character.
\end{corollary}

\begin{proof}
The witnesses have the form \(z c z c^{-1}\) with \(z\in M\),
so the lemma adds \(B_w=1\) to their previously proved
properties.  Their evaluations prove the asserted
injectivity exactly as before.

For the square abelian cover \(F_d\), every nontrivial
character line has scalar one.  The deck-trivial subspace
over \(\mathbf Q_r\) is the transfer of \(H_1(F,\mathbf Q_r)\),
on which multiplier one acts as the identity.  Thus the
whole \(H_1(F_d,\mathbf Q_r)\) is fixed.  Its integral
\(\mathbf Z_r\)-lattice is torsion-free, so it too is fixed.
Any finite abelian cover is dominated by one of these
square covers.  Transfer is injective over \(\mathbf Q_r\)
and is equivariant, so its homology is likewise fixed,
first rationally and then integrally.  This proves the
assertion for every coefficient prime.

Alternatively, on the image of \(B:M\to\Lambda^\times\),
complex conjugation acts by inversion.  An invariant
\(\mathbf F_p\)-character of this image is consequently
its own negative and is zero for odd \(p\).
The witness statement also excludes obtaining a nonzero
arithmetic target by jointly using these homology actions
with the earlier Kummer, word, and CM realizations.
It imposes no restriction on a genuinely nonmetabelian
mixed finite quotient.
\end{proof}
